\section{Complete finite-length RS certificate}
\label{app:rs-certificate}

This appendix proves the finite reductions used in
Theorem~\ref{thm:rs-finite} and specifies their executable verification.
Fix $m=16$, $\sigma\in\{4,8\}$, $h_0=m-\sigma$, $d=h_0+1$, and
$n=mB=mk/\sigma$, where $k=2^{20}$.  Set $\delta=3/5$ for $\sigma=4$
and $\delta=3/10$ for $\sigma=8$, and put $L=\lfloor\delta n\rfloor$.
All logarithms in the proof are natural.  We use
$Q=2^{127}-1$ as a lower bound on $q-1$.

An exact nonzero outer support meeting $t$ blocks has size $s$ and shortening
dimension $r$ satisfying
\[
 dt\le s\le mt,\qquad r=s-h_0t,\qquad 1\le r\le\sigma t.
\]
Its multiplicity is
\begin{equation}
 N_{t,s}=\binom Bt[z^s]
     \left(\sum_{\ell=d}^m\binom m\ell z^\ell\right)^t.
 \label{eq:rs-support-count}
\end{equation}
Every listed RS support is possible: a polynomial vanishing at its
complement has degree $m-\ell\le\sigma-1$ and is nonzero on the support.
Supports containing a nonempty local intersection smaller than $d$ have
empty exact-support events.  We omit them before applying the union bound.
Write $\widehat\mu$ for this restriction of the master sum~\eqref{eq:3}.
It still bounds the failure probability.

\subsection{A uniform projective envelope}

For $q-1\ge2n$ and $r\le n$, writing $Q'=q-1$ gives
\[
 \frac{q^r-1}{q-1}
 \le (Q')^{r-1}(1+1/Q')^r
 \le \frac{(Q')^{r-1}}{1-r/Q'}
 \le 2(Q')^{r-1}.
\]
The middle inequality follows by bounding each binomial coefficient
$\binom rj\le r^j$ and extending the resulting geometric series.
Consequently, for all $q\ge2^{127}$,
\begin{equation}
 \chi_q(r,c)\le 2\min\{1,Q^{r-1-c}\}.
 \label{eq:rs-chi-envelope}
\end{equation}
Indeed, $\min\{1,u^e\}$ is constant for $e\ge0$ and decreases with
$u\ge1$ for $e<0$.  For direct evaluation we use the sharper fact that
$\chi_q(r,c)=1$ for $c<r$.  For $c\ge r$, its value is
\[
 \sum_{j=1}^r\binom rj (q-1)^{j-1-c},
\]
which decreases with $q$.  Thus evaluation at $q=2^{127}$ supplies a
uniform upper bound, even when that integer is only a lower bound on the
actual field size.

\subsection{An envelope for all sparse supports}

Put $T=\lfloor n/1000\rfloor$ and choose $z=1/10$ for $\sigma=4$ and
$z=1/5$ for $\sigma=8$.  Define
\[
 \rho=\delta(1+z)^2,\quad D=1-\delta-mT/n,\quad
 V=(\delta D z^2)^{-1},\quad
 K_s=\frac{s}{n-s+1}\frac{\rho}{1-\rho}.
\]
The verifier checks $0<\rho<1$, $D>0$, $K_s<1$, $K_sQ>1$, and $V<Q$
using the extrema over $d\le s\le mT$.

First suppose $b\le s$.  Summing the last binomial in
$\cG_n(H,j,b)$ by the hockey-stick identity gives
\[
 \sum_{j=1}^L\cG_n(H,j,b)
 \le\frac{\binom{H-b}b\binom nb\binom L{H-b}}{\binom nH}.
\]
A nonzero term requires $H\le L+b\le L+s$.  The inequalities
\[
 \frac{\binom L{H-b}}{\binom n{H-b}}\le\delta^{H-b},\qquad
 \frac{\binom nb\binom n{H-b}}{\binom nH}
 \le\binom Hb\left(\frac n{n-H+1}\right)^b
\]
therefore imply
\begin{equation}
 \sum_{j=1}^L\cG_n(H,j,b)
 \le\binom Hb^2 D^{-b}\delta^{H-b}\le\rho^H V^b.
 \label{eq:rs-second-envelope}
\end{equation}
The last inequality uses $\binom Hb\le(1+z)^H/z^b$.
For the first stage,
\[
 \cG_n(s,H,a)\le
 2^s\left(\frac{s}{n-s+1}\right)^{s-a}
 \binom{H-1}{s-a-1}.
\]
This follows from $\binom{s-a}a\le2^s$,
$\binom{n-H}a\le\binom na$, and
$\binom na/\binom ns\le[s/(n-s+1)]^{s-a}$.
Extending the positive series over $H$ and using
\[
 \sum_{H\ge u}\binom{H-1}{u-1}\rho^H
   =\left(\frac\rho{1-\rho}\right)^u
\]
now bounds the two-stage placement sum by $2^sK_s^{s-a}V^b$.

To include all cancellations, write $K=K_s$.  For $b<r$ set
$e=a+b-r+1$.  Then
\[
 K^{s-a}V^b=K^{h_0t+1}(KV)^bK^{-e}.
\]
Apply~\eqref{eq:rs-chi-envelope} and extend $e$ to all integers.
The sums for $e\le0$ and $e\ge1$ are $1/(1-K)$ and $1/(KQ-1)$.
For $r\le b\le s$, every term has surplus cancellations; extending both
geometric sums in $a$ and $b-r$ gives the following bound for all $b\le s$:
\begin{equation}
 2^{s+1}\left[
 K^{h_0t+1}\left(\frac1{1-K}+\frac1{KQ-1}\right)
       \sum_{b=0}^{r-1}(KV)^b
 +\frac{K^{h_0t}(KV)^r}{Q(1-1/(KQ))(1-V/Q)}
 \right].
 \label{eq:rs-low-b}
\end{equation}

For $b>s$, discard the output-weight restriction.  Vandermonde's identity
gives
\[
 \sum_{j=1}^n\cG_n(H,j,b)=\binom{H-b}b\le\binom nb,
 \qquad \sum_{H=1}^n\cG_n(s,H,a)=\binom{s-a}a.
\]
Since $\sum_a\binom{s-a}a\le2^s$, dropping the additional factor $Q^{-a}$
and using $\binom nb\le n^b$ yields the complete tail bound
\begin{equation}
 \frac{2^{s+1}n^{s+1}Q^{-h_0t-2}}{1-n/Q}.
 \label{eq:rs-high-b}
\end{equation}

To aggregate~\eqref{eq:rs-low-b}, set
\[
\begin{aligned}
 U_0&=\sum_{\ell=d}^m\binom m\ell,& M_0&=2^mU_0,\\
 C&=\frac{mn}{n-mT+1}\frac\rho{1-\rho},&
 K_{\max}&=CT/n,\\
 K_{\min}&=\frac d{n-d+1}\frac\rho{1-\rho},&
 W&=\max\{1,K_{\max}V\},\\
 J_0&=\frac3m M_0 C^{h_0}W^\sigma,&
 x(t)&=J_0(t/n)^{h_0-1}.
\end{aligned}
\]
There are at most $\binom BtU_0^t\le(3B/t)^tU_0^t$ supports touching
$t$ blocks; the last inequality uses $e<3$.  Also $K_s\le Ct/n$ and
$\sum_{b<r}(K_sV)^b\le\sigma tW^{\sigma t}$.  Hence their combined
$b\le s$ contribution is at most
\begin{equation}
 (c_2t^2+c_0)x(t)^t,
 \label{eq:rs-t-envelope}
\end{equation}
where
\[
 c_2=\frac{2C\sigma}{n}
   \left(\frac1{1-K_{\max}}+\frac1{K_{\min}Q-1}\right),\qquad
 c_0=\frac2{Q(1-1/(K_{\min}Q))(1-V/Q)}.
\]
Use the disjoint intervals $[2,3]$, $[4,7]$, $[8,15]$, $[16,31]$,
$[32,63]$, and $[64,T]$.  On $[A,E]$, replace $x(t)$ by $x(E)<1$ and
extend the sum to all $t\ge A$.  The resulting geometric sums are rational:
\[
\begin{aligned}
 \sum_{t\ge A}x^t&=\frac{x^A}{1-x},\\
 \sum_{t\ge A}t^2x^t
 &=x^A\left(\frac{A^2}{1-x}
      +\frac{2Ax}{(1-x)^2}+\frac{x(1+x)}{(1-x)^3}\right).
\end{aligned}
\]
This covers every $2\le t\le T$, including growing sublinear block counts.
Finally, summing~\eqref{eq:rs-high-b} over all $1\le t\le T$ gives at most
\begin{equation}
 \frac{2n}{Q^2(1-n/Q)}\frac Y{1-Y},\qquad
 Y=\frac{3BM_0n^m}{Q^{h_0}}<1.
 \label{eq:rs-global-high-b}
\end{equation}
Here $s\le mt$ and $1/t\le1$ justify each relaxation.  The verifier checks
all positivity and convergence conditions by exact rational comparisons.

\subsection{Direct one-block terms and their intermediate tail}

For $t=1$, we retain the exact multiplicities~\eqref{eq:rs-support-count}.
For $b<r$, the final-weight sum has the finite identity
\begin{equation}
 \sum_{j=1}^L\cG_n(H,j,b)
 =\frac{\binom{H-b}b}{\binom nH}
   \sum_{v=0}^b\binom L{H-v}\binom{n-L}v.
 \label{eq:rs-final-cdf}
\end{equation}
Indeed, an $H$-subset has its $(H-b)$th order statistic at most $L$
exactly when it has at most $b$ elements above $L$.  Counting first by the
order statistic and then by that number of elements proves the identity.
We evaluate every $H\le512$, $a\le\lfloor s/2\rfloor$, and $b<r$
using~\eqref{eq:rs-final-cdf} and the exact first kernel.

For the remaining $H$, put $\delta_0=L/n$ and $D_s=1-\delta_0-s/n$.
Keeping the first binomial coefficient in the preceding envelope gives
\[
 \cG_n(s,H,a)\sum_{j\le L}\cG_n(H,j,b)
 \le A_{s,a,b}\binom{H-1}{s-a-1}\binom H{2b}\delta_0^H,
\]
where
\[
 A_{s,a,b}=
 \frac{\binom{s-a}a\binom na}{\binom ns}
 \frac{\binom{2b}b}{(D_s\delta_0)^b}.
\]
We used $\binom{H-b}b\binom Hb=\binom{2b}b\binom H{2b}$.
For an integer $w\ge1$ and $0<u<1$, define the rational function
\begin{equation}
 R_{w,b}(u)=
 \sum_{i=0}^{\min(w,2b)}
 \frac{w}{w+2b-i}
 \frac{(w+2b-i)!}{i!(w-i)!(2b-i)!}
 \left(\frac u{1-u}\right)^{w+2b-i}.
 \label{eq:rs-moment}
\end{equation}
It equals $\sum_{H\ge1}\binom{H-1}{w-1}\binom H{2b}u^H$.
To prove this, use $\binom{H-1}{w-1}=w\binom Hw/H$, count a pair of
subsets by its intersection size $i$, and sum the resulting union size
$N=w+2b-i$ with $\binom HN/H=\binom{H-1}{N-1}/N$.
Taking $u=(1+\delta_0)/2$ bounds every omitted $H>512$ by
\begin{equation}
 A_{s,a,b}(\delta_0/u)^{513}R_{s-a,b}(u).
 \label{eq:rs-H-tail}
\end{equation}
The script multiplies this expression by the projective factor and sums
every $a$ and $b<r$.  For $r\le b\le s$, it uses the second term
of~\eqref{eq:rs-low-b}, with the support's own $D_s$ and $K_s$.
Equation~\eqref{eq:rs-global-high-b} already covers $b>s$.

\subsection{Every remaining support type}

For $t>T$, let $k_0,k_d,\ldots,k_m$ count blocks of each local support size,
and set $p_\ell=k_\ell/B$.  The number of supports of that type is at most
$\exp(n\cS_{m,\sigma}(p))$ by the multinomial entropy inequality.
Normalize $s,H,a,b,j$ by $n$, writing $\alpha,\eta,x,y,\zeta$, and put
$\beta=\sum_{\ell=d}^m(\ell-h_0)p_\ell/m$.  Define
$f(A,B)=AH(B/A)$ on $0\le B\le A$, with $f(0,0)=0$.
Let $J_1$ be the sum of
\[
\begin{gathered}
 f(\alpha-x,x),\quad f(1-\eta,x),\quad f(\eta,\alpha-x),\\
 f(\eta-y,y),\quad f(1-\zeta,y),\quad f(\zeta,\eta-y).
\end{gathered}
\]
The finite exponent used by the verifier is
\[
 F=\cS_{m,\sigma}(p)+J_1-H(\alpha)-H(\eta)
       -\log Q\max\{0,x+y-\beta\}.
\]
The feasible region requires that $p$ be a probability vector,
$\alpha=\sum\ell p_\ell/m$, $p_0\le1-m(T+1)/n$, $0\le\zeta\le\delta$,
and that all six perspective terms have $0\le B\le A$.

Use $\binom vu\le\exp(vH(u/v))$ in the six numerators, after replacing
$\binom{H-1}{s-a-1}$ by $\binom H{s-a}$ and its second-stage analogue.
The lower bound $\binom nu\ge\exp(nH(u/n))/(n+1)$ controls both denominators.
Equation~\eqref{eq:rs-chi-envelope}, with its extra minus one discarded,
therefore bounds each support-type/trace contribution by
$2(n+1)^2\exp(nF)$.  There are at most $(B+1)^\sigma$ support types and
$(n+1)^4$ trace tuples.  Thus any uniform upper bound $F\le U$ gives
\begin{equation}
 \widehat\mu_{t>T}\le
       2(B+1)^\sigma(n+1)^6\exp(nU).
 \label{eq:rs-bulk-finite}
\end{equation}

Only $-H(\alpha)$ and $-H(\eta)$ are nonconcave terms of $F$.
The verifier covers
\[
 [d(T+1)/n,1]\times[d(T+1)/(2n),1]
\]
by rectangles with rational boundaries, and replaces these two terms by
their upper secants.  Every other entropy term has a global linear upper
bound.  For any $0<\theta<1$ and strictly positive probability vector $p^*$,
\[
 f(A,B)\le-B\log\theta-(A-B)\log(1-\theta),\qquad
 H(p)\le-\sum_\ell p_\ell\log p^*_\ell.
\]
Introduce $c\ge0$, $c\ge x+y-\beta$, with objective term $-c\log Q$.
The choice $c=\max\{0,x+y-\beta\}$ lies in $[0,1]$, since $x,y\le1/2$.
Each rectangle consequently has a linear upper objective $c_*+v_*\cdot v$
over rational constraints $Av\le b$, $Ev=e$, and $l\le v\le u$.
For any rational multipliers $\mu\ge0$ and $\nu$, put
$r_*=v_*-A^T\mu-E^T\nu$.  Then
\begin{equation}
 F\le c_*+\mu\cdot b+\nu\cdot e
       +\sum_i\bigl((r_*)_i l_i+\max\{0,(r_*)_i\}(u_i-l_i)\bigr).
 \label{eq:rs-dual-check}
\end{equation}
Arb intervals enclose the logarithms and residuals in this expression.
The numerical optimizer only proposes $p^*$, $\theta$, and multipliers;
their feasibility or optimality is not assumed.  The verifier normalizes
$p^*$ exactly and checks all required signs.

The covering is encoded as a binary tree of exact midpoint splits.  The
verifier reconstructs every boundary and rejects missing or extraneous nodes.
An infeasible leaf is accepted only if an exact necessary condition fails:
$\alpha_{\rm upper}\ge d(T+1)/n$,
$\alpha_{\rm lower}\le2\eta_{\rm upper}$, or
$\eta_{\rm lower}\le2\delta$.  The saved coverings have $21$ and $19$ leaves
for $\sigma=4$ and $\sigma=8$, respectively.

\subsection{Numerical result and reproduction}

Table~\ref{tab:rs-certificate} lists separate bounds for every region.
The complete verifier adds their interval enclosures before checking the
reported failure bound.  It obtains
\[
 \log_2\widehat\mu < -209.6008601171\quad(\sigma=4),\qquad
 \log_2\widehat\mu < -137.5682486336\quad(\sigma=8).
\]
These inequalities imply Theorem~\ref{thm:rs-finite}.

\begin{table}[ht]
\centering
\small
\begin{tabular}{@{}l@{\qquad}r@{\quad}r@{}}
\toprule
Region & $\sigma=4$ & $\sigma=8$\\
\midrule
$t=1$, $H\le512$, $b<r$ & $209.6008$ & $137.5682$\\
$t=1$, $H>512$, $b<r$ & $404$ & $690$\\
$t=1$, $r\le b\le s$ & $296$ & $243$\\
$2\le t\le T$, $b\le s$ & $250$ & $159$\\
$1\le t\le T$, $b>s$ & $1357$ & $862$\\
$t>T$, all traces & $1078$ & $5634$\\
\midrule
Complete bound & $209.60$ & $137.56$\\
\bottomrule
\end{tabular}
\caption{Certified lower bounds on $-\log_2$ of each contribution.
Every region includes all feasible first return counts and final weights.
The lengths and distances are those of Theorem~\ref{thm:rs-finite}.}
\label{tab:rs-certificate}
\end{table}

The exact directly summed terms with $t=1$, $H\le512$, and $a+b<r$ also
give a lower bound on $\widehat\mu$, because their projective factor is one.
The ratio of the complete upper bound to this partial sum is less than
$1+5\cdot10^{-13}$ for $\sigma=4$ and $1+3.3\cdot10^{-7}$ for $\sigma=8$.
Further numerical tightening of this same sum can gain less than a
millionth of a bit.  This comparison does not lower-bound the true failure
probability, since $\widehat\mu$ is itself a union bound.

The certificate artifact is provided in the accompanying
repository~\cite{chosenbaaartifact} and uses Python with
\texttt{python-flint}~0.9.0.  From the repository root, the complete
verification commands are
\begin{verbatim}
python -m pip install -r rs-certification/requirements.txt
python rs-certification/verify_complete.py
\end{verbatim}
The certificate directory contains both rational certificate proposals.
Verification needs neither NumPy nor SciPy.  Optional generation and
small-instance tests have separate dependencies and commands in that
directory's README.  Tests check the order-statistic identity, generating
series, intermediate tail, finite prefactors, and rejection of malformed
coverings.  These checks support the written reductions; the new analytic
and numerical certificate is not a Lean proof.
